problem:  
Let \(M = G/H\) be a compact, simply connected **generalized Wallach space**, i.e. a homogeneous space whose isotropy representation decomposes as \(\mathfrak{m} = \mathfrak{m}_1 \oplus \mathfrak{m}_2 \oplus \mathfrak{m}_3\) into three irreducible pairwise inequivalent \(H\)-submodules satisfying \([\mathfrak{m}_i, \mathfrak{m}_i] \subset \mathfrak{h}\) for each \(i\). Each \(G\)-invariant metric is determined by three positive parameters \(x_i\) giving \(g = x_1 Q|_{\mathfrak{m}_1} + x_2 Q|_{\mathfrak{m}_2} + x_3 Q|_{\mathfrak{m}_3}\), where \(Q\) is minus the Killing form on \(\mathfrak{g}\).  

Define \(P \subset \mathbb{R}^3_{>0}\) to be the region of parameters for which \(g(x_1,x_2,x_3)\) has strictly positive sectional curvature, and let \(E \subset \mathbb{R}^3_{>0}\) denote the algebraic locus of \(G\)-invariant Einstein metrics, given by the standard homogeneous Einstein equations.  

**Question:**  
Is it true that for every generalized Wallach space \(G/H\) with \(P \neq \emptyset\), the closure \(\overline{P}\) in \(\mathbb{R}^3_{>0}\) intersects \(E\)?  
Equivalently, can one always continuously deform a positively curved invariant metric toward an Einstein metric within the invariant metric moduli, or do there exist generalized Wallach spaces for which \(P\) and \(E\) are separated by a region of nonpositive curvature in parameter space?  

Why is it a "good" problem:  
This formulation is mathematically precise, logically valid, and targets a subtle and unexplored interaction between **positive sectional curvature** and the **Einstein property** within the finite-dimensional moduli space of invariant metrics on generalized Wallach spaces—a class that includes the Wallach and Aloff–Wallach families. The question proposes a concrete geometric condition (intersection of regions in parameter space) representing an intrinsic rigidity or continuity phenomenon: does the corner of positivity always “touch” the Einstein locus?

Resolving this would yield new insights bridging **homogeneous Einstein geometry**, **curvature inequalities**, and **Lie-theoretic structure constants**, with immediate implications for the broader classification of positively curved homogeneous spaces. The statement remains elementary to express (“Do the positive curvature metrics accumulate on the Einstein ones?”) yet demands sophisticated analytic and algebraic methods to answer, embodying the principle of “simple to state, deep to solve” that typifies a genuinely good research problem.